\subsection{Fiducialization Efficiency}
In order to reduce the impact of background events originating from the detector materials and the surrounding environment, a fiducial volume cut is typically applied to the reconstructed neutrino events in low energy--analyses. This cut is based on the reconstructed position of the primary \ac{TPC} cluster associated with the neutrino interaction. The fiducial volume is defined as a region within the detector that is at least x~\unit{\centi\meter} away from any external detector boundary. This distance can be chosen in any analysis to optimize the signal-to-background ratio.

\begin{figure}[ht!]
    \centering
    \includegraphics[width=0.49\linewidth]{Images/hd_1x2x6_lateralAPA_marley_official_Vertex_FiducialEfficiency.png}
    \includegraphics[width=0.49\linewidth]{Images/vd_1x8x14_3view_30deg_nominal_marley_official_Vertex_FiducialEfficiency.png}
    \caption{Fiducialization efficiency as a function of the distance to the detector boundary for (left) \ac{HD} and (right) \ac{VD}.\label{fig:FIDUCIALIZATION_EFFICIENCY}}
\end{figure}